\documentclass[10pt,a4paper]{amsart}
\usepackage{amsmath,amssymb,amsthm,graphicx}
\usepackage{verbatim}
\newenvironment{reference}{\comment}{\endcomment}
%\newenvironment{reference}{}{}
\newenvironment{slogan}{\comment}{\endcomment}
\newenvironment{history}{\comment}{\endcomment}
% For commutative diagrams we use Xy-pic
\usepackage[all]{xy}

% We use 2cell for 2-commutative diagrams.
\xyoption{2cell}
% Theorem environments.
%
\theoremstyle{plain}
\newtheorem{theorem}[subsection]{Theorem}
\newtheorem{proposition}[subsection]{Proposition}
\newtheorem{lemma}[subsection]{Lemma}

\theoremstyle{definition}
\newtheorem{definition}[subsection]{Definition}
\newtheorem{example}[subsection]{Example}
\newtheorem{exercise}[subsection]{Exercise}
\newtheorem{situation}[subsection]{Situation}

\theoremstyle{remark}
\newtheorem{remark}[subsection]{Remark}
\newtheorem{remarks}[subsection]{Remarks}

\numberwithin{equation}{subsection}

% Macros
%
\def\lim{\mathop{\mathrm{lim}}\nolimits}
\def\colim{\mathop{\mathrm{colim}}\nolimits}
\def\Spec{\mathop{\mathrm{Spec}}}
\def\Hom{\mathop{\mathrm{Hom}}\nolimits}
\def\Ext{\mathop{\mathrm{Ext}}\nolimits}
\def\SheafHom{\mathop{\mathcal{H}\!\mathit{om}}\nolimits}
\def\SheafExt{\mathop{\mathcal{E}\!\mathit{xt}}\nolimits}
\def\Sch{\mathit{Sch}}
\def\Mor{\mathop{\mathrm{Mor}}\nolimits}
\def\Ob{\mathop{\mathrm{Ob}}\nolimits}
\def\Sh{\mathop{\mathit{Sh}}\nolimits}
\def\NL{\mathop{N\!L}\nolimits}
\def\CH{\mathop{\mathrm{CH}}\nolimits}
\def\proetale{{pro\text{-}\acute{e}tale}}
\def\etale{{\acute{e}tale}}
\def\QCoh{\mathit{QCoh}}
\def\Ker{\mathop{\mathrm{Ker}}}
\def\Im{\mathop{\mathrm{Im}}}
\def\Coker{\mathop{\mathrm{Coker}}}
\def\Coim{\mathop{\mathrm{Coim}}}

% Boxtimes
%
\DeclareMathSymbol{\boxtimes}{\mathbin}{AMSa}{"02}

%
% Macros for moduli stacks/spaces
%
\def\QCohstack{\mathcal{QC}\!\mathit{oh}}
\def\Cohstack{\mathcal{C}\!\mathit{oh}}
\def\Spacesstack{\mathcal{S}\!\mathit{paces}}
\def\Quotfunctor{\mathrm{Quot}}
\def\Hilbfunctor{\mathrm{Hilb}}
\def\Curvesstack{\mathcal{C}\!\mathit{urves}}
\def\Polarizedstack{\mathcal{P}\!\mathit{olarized}}
\def\Complexesstack{\mathcal{C}\!\mathit{omplexes}}
% \Pic is the operator that assigns to X its picard group, usage \Pic(X)
% \Picardstack_{X/B} denotes the Picard stack of X over B
% \Picardfunctor_{X/B} denotes the Picard functor of X over B
\def\Pic{\mathop{\mathrm{Pic}}\nolimits}
\def\Picardstack{\mathcal{P}\!\mathit{ic}}
\def\Picardfunctor{\mathrm{Pic}}
\def\Deformationcategory{\mathcal{D}\!\mathit{ef}}
\usepackage{cleveref}
\crefname{originallemma}{original lemma}{original lemmas}
\Crefname{originallemma}{Original lemma}{Original lemmas}
\makeatletter
\newlabel{lemma-representable-decent}{{lemma-representable-decent}{source}{}{sourcepointer-1}{}}
\newlabel{lemma-representable-decent@cref}{{[originallemma][0][]lemma-representable-decent}{[1][0][]source}{}{}{}}
\makeatother
\crefname{originallemma}{original lemma}{original lemmas}
\Crefname{originallemma}{Original lemma}{Original lemmas}
\makeatletter
\newlabel{lemma-specialization}{{lemma-specialization}{source}{}{sourcepointer-2}{}}
\newlabel{lemma-specialization@cref}{{[originallemma][0][]lemma-specialization}{[1][0][]source}{}{}{}}
\makeatother
\crefname{originalreference}{original reference}{original references}
\Crefname{originalreference}{Original reference}{Original references}
\makeatletter
\newlabel{topology-definition-generic-point}{{topology-definition-generic-point}{source}{}{sourcepointer-3}{}}
\newlabel{topology-definition-generic-point@cref}{{[originalreference][0][]topology-definition-generic-point}{[1][0][]source}{}{}{}}
\makeatother
\title{The topological space of a decent algebraic stack is Kolmogorov}
\author{The Stacks Project Authors}
\date{stacks-project:stacks-morphisms.tex:local-snapshot-20261012; GNU FDL 1.2 or later}
\begin{document}
\maketitle
\noindent\textit{Editorial note.} All delivered mathematical text is the author's own, copied from the pinned Stacks Project chapter source sources/stacks-algebraic-stacks-morphisms/pinned-original-source.tex; no statement or proof step was rewritten, reordered or removed.\par\medskip
\noindent\textit{Editorial note.} Preamble provenance: the Stacks Project keeps one single shared preamble for every chapter; this card takes it directly from the staged sources/stacks-algebraic-stacks-morphisms/preamble.tex as three verbatim spans (lines 8-12, 14-18 and 154-217), so no derived file was written and nothing was added or redefined. The wrapper class and the base packages are supplied by the assembler.\par\medskip
\noindent\textit{Editorial note.} Bibliography provenance: neither the counted statement nor its proof contains a citation (this lane rescans the whole span before writing the card), so the card carries no bibliography and no entry was reconstructed.\par\medskip
\noindent\textit{Editorial note.} Selection provenance: this card was opened under the file-based control contract (CONTROL.md), with the mechanical pre-card gate run on its statement and proof interval before the card was written.\par\medskip
\noindent\textit{Allowed standard prerequisites (editorial).}\begin{itemize}
\item Decency of the reduced closed substack supported on the closure of a point. Invoked inside the counted proof (source lines 10528-10530). Allows the reduction to the irreducible case inside a decent closed substack.
\item The specialization lifting lemma for smooth morphisms from schemes to stacks. Invoked inside the counted proof (source lines 10549-10552). Lifts a specialization between generic points to a specialization in the presenting scheme.
\item The characterization of points through generic points of a smooth presentation. Used inside the counted proof (source lines 10537-10543). Produces a third generic point of the space from a generic point of a component of the presentation.
\end{itemize}\medskip
\ifdefined\hypertarget\hypertarget{sourcepointer-1}{}\fi
\noindent\textit{Source reference.} Original lemma lemma-representable-decent: sources/stacks-algebraic-stacks-morphisms/pinned-original-source.tex, source lines 10531--10531. Navigation pointer only: the counted text refers to the Stacks item lemma-representable-decent, which belongs to another chapter of the project and is not part of this card. The pointer carries the target's own label as its number because this chapter's staged snapshot does not contain the other chapter's tag table; its content is not reproduced here.\par\medskip
\ifdefined\hypertarget\hypertarget{sourcepointer-2}{}\fi
\noindent\textit{Source reference.} Original lemma lemma-specialization: sources/stacks-algebraic-stacks-morphisms/pinned-original-source.tex, source lines 10553--10553. Navigation pointer only: the counted text refers to the Stacks item lemma-specialization, which belongs to another chapter of the project and is not part of this card. The pointer carries the target's own label as its number because this chapter's staged snapshot does not contain the other chapter's tag table; its content is not reproduced here.\par\medskip
\ifdefined\hypertarget\hypertarget{sourcepointer-3}{}\fi
\noindent\textit{Source reference.} Original reference topology-definition-generic-point: sources/stacks-algebraic-stacks-morphisms/pinned-original-source.tex, source lines 2283--2283. Navigation pointer only: the counted text refers to the Stacks item topology-definition-generic-point, which belongs to another chapter of the project and is not part of this card. The pointer carries the target's own label as its number because this chapter's staged snapshot does not contain the other chapter's tag table; its content is not reproduced here.\par\medskip
\begin{lemma}
\label{lemma-kolmogorov}
Let $\mathcal{X}$ be a decent algebraic stack.
Then $|\mathcal{X}|$ is Kolmogorov (see
Topology, Definition \ref{topology-definition-generic-point}).
\end{lemma}

\begin{proof}
Let $x_1, x_2 \in |\mathcal{X}|$ with $x_1 \leadsto x_2$ and
$x_2 \leadsto x_1$. We have to show that $x_1 = x_2$.
Let $\mathcal{Z} \subset \mathcal{X}$ be the reduced closed
substack with $|\mathcal{Z}|$ equal to
$\overline{\{x_1\}} = \overline{\{x_2\}}$.
By Lemma \ref{lemma-representable-decent}
we see that $\mathcal{Z}$ is decent.
After replacing $\mathcal{X}$ by $\mathcal{Z}$ we reduce to
the case discussed in the next paragraph.

\medskip\noindent
Assume $|\mathcal{X}|$ is irreducible with generic points
$x_1$ and $x_2$. Pick an affine scheme $U$ and $u_1, u_2 \in U$
and a smooth morphism $f : U \to \mathcal{X}$ such that
$f(u_i) = x_i$. Then we find a third point $u_3 \in U$
which is the generic point of an irreducible component of $U$
whose image $x_3 \in |\mathcal{X}|$ is {\it also} a
generic point of $|\mathcal{X}|$. Namely, we can simply
choose $u_3$ any generic point of an irreducible component
passing through $u_1$ (or $u_2$ if you like).
In the next paragraph we will show that $x_1 = x_3$ and $x_2 = x_3$
which will prove what we want.

\medskip\noindent
By symmetry it suffices to prove that $x_1 = x_3$. Since $x_1$
is a generic point of $|\mathcal{X}|$ we have a specialization
$x_1 \leadsto x_3$.
By Lemma \ref{lemma-specialization} we can find a specialization
$u'_1 \leadsto u_3$ in $U$ (!) mapping to $x_1 \leadsto x_3$.
However, $u_3$ is the generic point of an irreducible component
and hence $u'_1 = u_3$ as desired.
\end{proof}

\end{document}
